\section{Intrinsic Reward Dynamics}
\label{sec:results_intrinsic_reward_dynamics}

The intrinsic reward is intended to shape exploration during training rather than to define the final task objective.
For this reason, the evaluation protocol in Section~\ref{sec:methodology_evaluation_methodology} reports only extrinsic return, while the training diagnostics examine how much intrinsic reward is added to the PPO update.
Figure~\ref{fig:results_reward_decomposition} decomposes rollout rewards into extrinsic and intrinsic components.

\begin{figure}[!hbp]
\centering
\includegraphics[width=\glpefigwidth,height=\glpefigheight,keepaspectratio]{chapters/05_results_and_discussion/figures/train-reward-decomp-baselines.png}
\caption{Training-time rollout reward decomposition. Solid lines show mean extrinsic reward per environment step collected during PPO rollouts. Dashed lines show the intrinsic shaping component after scaling and clipping. Evaluation metrics reported in Sections~\ref{sec:results_learning_curves}--\ref{sec:results_wall_clock_efficiency} use extrinsic return only.}
\label{fig:results_reward_decomposition}
\end{figure}

The diagnostic curves show that GLPE and GLPE without gating apply intrinsic shaping primarily as a transient auxiliary signal.
The intrinsic coefficient is tapered during training, so the contribution of intrinsic reward decreases as the policy approaches the later stages of the interaction budget.
This schedule reduces the chance that the final policy is optimized mainly for auxiliary novelty or progress rather than environment return.

The reward decomposition also clarifies why training rewards cannot be compared directly with evaluation returns.
During training, the policy update uses the sum of extrinsic reward and scaled intrinsic reward.
During evaluation, only the environment return is measured.
Consequently, a large intrinsic component can accelerate learning if it guides useful exploration, but it can also degrade performance if it encourages behavior that is not aligned with extrinsic return.

For prediction-error methods, intrinsic reward may remain large in regions that are difficult to predict, even when those regions no longer provide useful learning opportunities.
GLPE addresses this issue through local learning-progress estimation and optional gating.
The observed tapering and gate suppression are therefore complementary mechanisms: tapering reduces the global strength of intrinsic shaping over time, while gating suppresses specific regions that appear unproductive before the global coefficient has decayed.

\FloatBarrier